\section{怎样进行构思}
审题之后，就要根据题目限制的范围、要求和体裁，以及题目没加限制的领域，开始打腹稿，把自己已有的材料认识，经过一番精细的思考和联想，组织成一篇文章的雏形，这个过程就叫构思。这如同盖房子要设计蓝图一样，究竟要盖什么样的房子，哪些材料可用，哪些材料不可用，哪儿是房子的主体，哪儿是副件，都要在动工之前设计好。不然的话，所盖的房子轻则简陋难看，杂乱无章，重则东倒西歪，需要推倒重来。构思就是为要写的文章设计图样，这是写每一篇文章必不可少的过程。如果看了题目就写，而不对文章的全貌作一番精细的思索，这只能是盲目的写作，写出的文章十之八九是主题不明，缺乏条理，结构松散，内容空洞。

构思既然这么重要，我们在下笔之前，怎样才能做好这件事呢？

\textbf{1. 根据题目，充分展开联想和想象，准备作文的素材。}

一个作文题目只定下大致范围，不曾限定的范围还很开阔，可写的内容也很多；就是在所定的范围内，怎样选取材料，怎样组织材料，怎样突出中心思想，也是各式各样的。因此，作文的内容和写法有广阔的天地，只待我们广开思路去寻求，去发现。

广开思路的基本途径就是展开联想和想象。有人说，听、说、读、写这四种语文能力的总枢纽就是一个“想”字，因为不善于动脑筋，想问题，是很难培养出听、说、读、写的语文能力的。我们面对着一个个作文题，要立即开动思想机器，进行充分的回忆，多方的思考，尽量多想些与题目有关的人和事，想得越具体越细致越好。要由此及彼、由己及人、由今及古，由近及远、由小及大、由表及里地想，反复不断地开拓下去，这就是联想；对不在眼前的事物想出它的具体形象，把脑子里所留下的印象组合起来，化为具体的可捉摸的东西，这就是想象。联想和想象不是玄妙的东西，而是人人都可以具有的能力，只不过有的人强，有的人弱罢了。写作文正是锻炼联想和想象能力的好机会，如果没有一点联想和想象的能力，是不可能进行写作这种创造性的劳动的。

具有各种不同联系的事物，反映在人们的头脑中，会形成各种不同的联想：在空间或时间上接近的事物形成接近联想，如由面前的水库想到将来可在这个水库的一侧建立水电站；有相似特点的事物形成类似联想，如由蜻蜓想到飞机，由江青想到吕后；有对立关系的事物形成对比联想，如由今天的大好形势想到十年动乱；有因果关系的事物形成因果联想，如由良好学风的树立，想到学习成绩的普遍提高，由丰收景象想到农民的辛勤劳动。在作文中，充分展开联想，就能打开思路。

例如，我们面对着《伟人事迹给我的启示》这个题目，首先就得想一想，什么样的人才够称得上“伟人”，有哪些门类（如政治、文学、艺术、科学、军事等）的伟人，他们的代表人物有哪些，古代的有哪些，现代的有哪些，当代的有哪些，他们的典型事迹各是什么，他们的事迹有哪些性质相同或相似的地方，他们这方面性质相同的事迹能够给我什么样的启示，他们那方面性质相同的事迹能够给我什么样的启示，有哪些由于受到伟人事迹的启示而得出的名言警句，等等。这样一想，可供选用的素材就开始多起来了。这种分门别类、因类联想的方式可以叫做“横向联想法”吧。
\newpage

我们还可以采取由人想事、由事想言、由言想其精神境界、由其精神境界想到给自己的启示的“纵向联想法”，从伟人中选取最富有代表性的三五个，集中地深入地想下去，想想他们成为伟人的关键点在哪里，想想他们所有事迹中最足以流传后世的东西是什么，想想最能给自己以教育、以鼓舞、以鞭策的启示又是什么。例如我国文学史上第一位伟大的文学家屈原，为什么能永垂不朽、受到中国人民子孙万代的敬仰和爱戴呢？就在于他忠于祖国，他那“九死其犹未悔”的执着的爱国精神使他永标青史，感召华夏子孙，这就给我们以启示：忠于祖国的人是不死的，将永远活在人民的心中，我们就该做这样的人。再如伟大的无产阶级革命战士方志敏之所以永垂不朽，就在于他忠于无产阶级革命事业，他身陷囹圄，威武不屈，并用敌人叫他写“自白书”的纸张，挥笔疾书《可爱的中国》、《狱中纪实》等革命雄文，来教育革命后代，表现了对祖国对人民对党的无比忠诚，这就是他伟大之处，他的事迹启示我们：忠于革命，宁死不屈的无产阶级战士是不死的，将永远活在人民的心中，我们就该做这样的人。其他如伟大的民族英雄岳飞、文天祥、伟大的史学家司马迁、伟大的文学家鲁迅、伟大的国际主义战士罗盛教、邱少云、等等，都可以作类似的联想。这样一想，对于可供选用的素材理解得更进一步了，就可以为下一步去粗存精的工作打好基础。

联想和想象，只要紧扣所写的题目，可以向前想，思接千古；向后想，想到未来子孙，向大处想，可包揽宇宙；向细处想，可明辨毫芒。这在记叙文写作中尤其重要，因为记叙文必须内容丰富，情节生动。我们可以打开记忆的闸门，先由想到的一点突出的事，纵向地扩展开去，逼紧我们的想象，事情开始怎样，然后又怎样，接着又怎样······其中最动人的地方具体情形又怎样？也可以由一件事横向地扩展开去，想起了甲事，随着联想到乙事、丙事······，其中人物神态的捕捉，音容笑貌的描写，心理活动的刻画，场面的再现，环境的渲染，等等，更需要我们用心地想象。如果能想象到如同其声，如见其人的地步，那么写出来的文章一定是非常动人的。

\textbf{2. 分析和提炼素材，确定好中心}

缺失282页下部分

缺失282页下部分

缺失282页下部分

缺失282页下部分

\newpage
\noindent
当然只是作为举例而设想的）······如果认真分析一下，最后三种类型既不合乎题意，又没有积极意义，可以抛开。前面的几种类型，好象都既合题意，又有积极意义，就应该再加分析。比如汇报同学干好事，除了自己不遵约之外，倒是表扬做了好事而不留名的“雷锋精神”，这个中心意思与“净友”的含义有一定距离，也可以抛开，而这个素材可以保存起来以后用。其他几种类型，经过分析和提炼，则都能从不同方面表达“只有直言规劝自己上进的朋友才是真正的朋友”这个中心意思，这样的中心意思，既切题，又明确，而且有积极意义，写出的文章一般会合乎要求的。至于议论文的题目，往往已点出了中心意思，如《勤能补拙》、《学贵有恒》、《我们要为振兴中华而勤奋学习》、《一屋不扫，何以扫天下》、《千里之行，始于足下》等，同《净友》这类记叙文题目一样，只能在几种中心意思中确定写某一种中心意思。在并不能明显表明中心意思的作文题中，情况就有所不同了，如《我的第一位班主任》、《亲人》、《我遇到这样一件新鲜事》、《雨中》、《灯下》、《毕业前夕》、《说“韧”》、《荣辱辨》、《百米冲刺的启示》、《在起跑线上》、《花》，等等，它们可以在同一个题目中确立表现出各不相同的中心意思来；当然，同一个中心意思的文章，也可以用不同的题目去表示。了解这一点，作文时就可灵活处理了。

\textbf{3. 围绕中心意思，选择材料和组织材料。}

当我们确定了一个明确的、集中的、有积极意义的中心意思之后，就可以围绕这个中心意思，对开始构思时所想的那些素材进行选择，决定取舍了。凡是不能表现中心意思的素材，都应该毫不可惜地舍弃（如上述《净友》的后四条素材就是这样），凡是能表现中心意思的素材都留下来作进一步的思考，这是第一点。在能表现中心意思的素材中进行比较，把那些最有表现力、最有说服力的素材留下来，把那些表现力、说服力较差的素材或者舍弃掉，或者仅作文章的中心材料的陪衬、补充。如上述《净友》的前几条关于学习方面的素材，比较一下就会发现：如果用批评自己平时松松垮垮、一遇考试就临时抱佛脚来写《净友》，不是不可以，但是内容比较空洞，不容易感动人，说服人；而如果用批评自己不认真学习而去偷看电影或偷看小说作为作文的材料，则容易写得具体感人，能较好地表达《净友》的中心意思，因此应选取后者，不选取前者，这是第二点。在选材时，还要注意到材料的新颖深切，有自己独特的体验和亲切的感受。例如：自己上课时偷看小说，好友一批评就不看了。但是，那一天为了偷看电影《少林寺》，自己编造了一张病假条送给班主任；当同学们全神贯注地上晚自习时，自己却正为少林十三僧的精彩武打所陶醉；谁知看完电影满腹高兴地回到家中时，自己的好友杨军却一脸不愉快地等着自己，原来他听说我病了，特来看望我，而当他从我小妹那儿知道我是去偷看《少林寺》，他就火了，一直等着我回来；现在他看见我兴高彩烈的样儿，不由得噼噼啪啪地跟我打起了机关枪，指责我一不该违反学校纪律，二不该向老师撒谎，三不该撕毁我们共同攻关的协议；“哟嗬嗬，不就看一场电影嘛，罪名可不少呢······”我这一咕噜不打紧，把他惹得更加火冒三丈，说我错了还不认错，干脆两下拉倒；我这才慌起来，赶紧赔不是，他还不敢松，硬说知错认错还不行，必须要改错，不由分说地拉我坐下来，陪着我一起完成当天数学和物理作业，其中一道数学题在演算时搞错了一个数，他非叫我从头检查起不可，好容易答数算对了，时钟已敲了十下，我说今天就做到这儿吧，明天我再来补，他坚持要今日事今日毕，我只好继续跟他一起做下去······。我们把以上这两份素材比较一下，可以看出，前者一般化，后者则有深切的感受，有新颖独到的地方，当然应该选取后者而舍弃前者，这是第三点。

选好了材料，就要考虑怎样组织安排材料，使文章成为一个有机的整体。尽管怎样组织安排材料没有一成不变的框框，不过对于初学写作的人来说，了解一些常见的写法却是必要的。

例如：先写什么，后写什么？这是一个按什么顺序来组织安排材料的问题。一般说来，我们可以按照事件发展的过程、空间位置的变化，或者按照问题的性质分类，来安排组织材料。

按事件发展的过程来安排组织材料，可以顺着事物发展的先后次序一步步地写下去，这叫顺叙，如峻青的《党员登记表》；也可以把事件的结局提前来写，然后再从头说起，这就叫倒叙，如鲁迅的《祝福》；也可以在叙述过程中插进一段过去的回忆或补充的事实，然后再接着原来叙述的地方继续叙述下去，这就叫插叙，如鲁迅的《故乡》，写到“我”和“母亲”谈到闰土时，插入一段回忆少年闰土和“我”的纯真友谊的片断，这种插叙的方法多半用于记叙文。需要注意：在运用倒叙或插叙的方法时，一定要把起、迄的时间关系交代清楚，往往在起、迄的地方用一些过渡的话语加以表明，使读者清晰地了解事情的先后顺序、来龙去脉；倒叙或插叙部分仍要按时间顺序写，不要使时间关系发生混乱。

按空间位置的变化来安排组织材料，可以顺着地点、方位的不同一层层地写下去。这一种方法大多在写游记、散文和说明文时被采用。如朱自清的《威尼斯》、周定舫的《人民英雄永垂不朽》、黄传惕的《故宫博物院》、孙世恺的《雄伟的人民大会堂》等，就是这样的例子。

按问题的性质分类来安排组织材料，即把表现中心意思的材料，按照它们的性质加以分类，相同的归在一起，写到一个层次里，使几个内容不同的层次合起来能表达文章的中心思想。这一种方法多在写议论文时被采用。如毛主席的《民族的科学的大众的文化》，就是把题目化成为什么说这种新民主主义文化是民族的、是科学的、是大众的这样三个问题来组织材料的。

其实，这几种方法在具体写作时往往可以交织使用。我们应该努力学会综合运用这些方法，使文章组织安排得更灵活、更富有变化，更能鲜明突出地表达中心思想。

再如：什么地方要详写，什么地方要略写？不言而喻，能够突出表现中心意思的地方要详，否则要略；重点人物、重要场面的描写要详，次要人物、次要场面的描写要略；记叙文写的事件要详，议论文中所摆的事实要略；重点分析处要详，一般交代性文字要略。详写了次要的方面，会使中心思想不突出，以致喧宾夺主；完全不写次要的方面，又可能使情节不完整，前后不连贯，或者使文章显得单调，同样也不能很好地表达中心思想。所以，对那些刻画重点人物形象、能够充分表达中心思想而需要详写的地方要不惜笔墨，写得淋漓尽致；对那些只起衬托、渲染、过渡、说明作用的

不需要详写的次要部分则要惜墨如金，舍得割爱。这样就能做到有详有略，主次分明，而又互相映衬，突出中心。

以上所说的这些方法，我们在构思文章考虑材料的安排组织时应该注意到，并且努力灵活地恰当地加以运用。
\newpage